\documentclass[11pt]{article}
\usepackage[a4paper,margin=25mm]{geometry}
\usepackage{amsmath,amssymb,amsthm,mathtools,microtype,graphicx,needspace}
\usepackage[colorlinks=true,linkcolor=blue,citecolor=blue,urlcolor=blue]{hyperref}
\newtheorem{theorem}{Theorem}
\newtheorem{proposition}[theorem]{Proposition}
\newtheorem{lemma}[theorem]{Lemma}
\newtheorem{corollary}[theorem]{Corollary}
\theoremstyle{remark}\newtheorem{remark}[theorem]{Remark}
\newcommand{\D}{\mathbb D}
\newcommand{\T}{\mathbb T}
\newcommand{\R}{\mathbb R}
\newcommand{\C}{\mathbb C}
\newcommand{\Rea}{\operatorname{Re}}
\newcommand{\Ima}{\operatorname{Im}}
\newcommand{\avg}[1]{\left\langle #1\right\rangle}
\newcommand{\ee}{\mathrm e}
\title{A sharp polynomial Lyapunov threshold\\for a stable cubic family}
\author{Lluis Eriksson}
\date{21 October 2026 (version 2)}
\begin{document}
\maketitle
\begin{abstract}
We determine exactly when the globally asymptotically stable cubic
family introduced by Liu and Fitzsimmons admits polynomial Lyapunov
certificates. In their normalization, with damping \(a>0\) and angular
forcing \(b\in\R\), the criterion is
\[
 |b|^{2/3}<1+a^{2/3}.
\]
The same strict inequality characterizes positive homogeneous
polynomial certificates with weak nonpositive Lie derivative, strict
polynomial certificates, and local real-analytic certificates.
Equality admits none of them. Below the threshold we construct a
certificate as a finite product of positive quadratic forms.
The proof pairs positive harmonic functions with an explicit analytic
square-root test and constructs a complementary limiting harmonic
pair at the boundary. Poisson smoothing and equal-weight quantiles
then produce a polynomial with a strict margin.
For fixed \(b\), an explicit sufficient degree is
\(O_b(\delta^{-3})\), where \(\delta\) is the damping excess above the
threshold, while
the minimum degree diverges; no matching degree rate is asserted.
The threshold persists for every single even angular harmonic and
nonzero angular speed. Classical binary-form factorization gives
the corresponding sum-of-squares consequence without increasing the
Lyapunov degree. The family, its global stability, the original
counterexample to a polynomial converse theorem and the nonexistence of
weak polynomial certificates for \(|b|\ge2(a+1)\) are prior results
of Liu and Fitzsimmons.
Our classification concerns certificates, not a change in dynamical
stability.
\end{abstract}

\section{The family and the exact certificate region}
For \(a>0\) and \(b\in\R\), consider
\begin{equation}\label{eq:field}
 \begin{aligned}
 \dot x&=-a(x^2+y^2)x+b(x^2-y^2)x-(x^2+y^2)y,\\
 \dot y&=-a(x^2+y^2)y+b(x^2-y^2)y+(x^2+y^2)x .
 \end{aligned}
\end{equation}
In polar coordinates \(x=r\cos\theta,y=r\sin\theta\), this is
\begin{equation}\label{eq:polar}
 \dot r=r^3(-a+b\cos2\theta),\qquad \dot\theta=r^2 .
\end{equation}
Liu and Fitzsimmons introduced this family \cite[Sec.~IV-D]{liu}.
They prove, for every \(a>0\) and \(b\in\R\), global asymptotic
stability of the origin through the Lyapunov function
\(r^2\ee^{-b\sin2\theta}\), and, for \(a>0\) and \(|b|\ge2(a+1)\), that
no positive definite homogeneous polynomial has nonpositive Lie
derivative \cite[Theorem IV.3]{liu}. The bound behind the latter is
\(|b|/2<a+1\), obtained from the second Fourier coefficient of the
normalized logarithmic derivative together with \(|g_2|<1\)
\cite[Lemma III.3]{liu}. For the single field \((a,b)=(1,5)\)
they also exclude local real-analytic Lyapunov functions by a
leading-term argument \cite[Theorem II.2(4), Sec.~III-D]{liu}; the
family version \cite[Theorem IV.3]{liu} is stated for homogeneous
polynomials only, and the extension of the analytic statement to the
region \(|b|\ge2(a+1)\) is a routine consequence of their argument
(see Section~\ref{sec:local}), not a statement we find there. Their work
already refutes the homogeneous polynomial Lyapunov converse. We
determine the exact polynomial and local analytic certificate region of
this family; it lies strictly inside \(|b|<2(a+1)\) (Remark~\ref{rem:inside}).
The test functions below pair the whole sequence of Fourier
coefficients at once, not only the second one.

For a vector field \(f\), write \(L_fV=\nabla V\cdot f\).
A \emph{weak certificate} means \(V(0)=0\), \(V>0\) away from the
origin, and \(L_fV\le0\). A \emph{strict certificate} has
\(L_fV<0\) away from the origin. For a local analytic certificate
these conditions are required on a punctured neighborhood of zero.
All homogeneous certificates below have positive degree.

\Needspace{15\baselineskip}
\begin{theorem}\label{thm:main}
For \eqref{eq:field}, the following are equivalent:
\begin{enumerate}
\item \( |b|<(1+a^{2/3})^{3/2}\).
\item There exists a positive definite homogeneous polynomial weak
certificate.
\item There exists a strict homogeneous polynomial certificate that
is a finite product of positive definite real quadratic forms.
\item There exists a local real-analytic weak certificate.
\end{enumerate}
Thus equality at the threshold is excluded even for weak certificates.
Every system in \eqref{eq:field} is nevertheless globally
asymptotically stable.
\end{theorem}

The stability assertion is recalled from \cite[Sec.~IV-D]{liu}, not new.
For orientation, its nonpolynomial certificate is
\[
 H(r,\theta)=r^2\exp(-b\sin2\theta),\qquad
 L_fH=-2ar^2H.
\]
It is positive, coercive, continuously differentiable at the origin,
and smooth away from it. The new assertion in Theorem~\ref{thm:main}
is the exact algebraic and analytic threshold.

\begin{figure}[t]
\centering
\includegraphics[width=.78\textwidth]{images/phase-boundary.pdf}
\caption{Certificate regions for the stable cubic family. All displayed
parameters with \(a>0\) give globally asymptotically stable dynamics.
The solid curve is excluded from the polynomial and local analytic
certificate region. The dashed line marks the sufficient
nonexistence condition for polynomial certificates in \cite[Theorem IV.3]{liu};
it is not a stability boundary.}
\label{fig:phase}
\end{figure}

Unbounded required polynomial degree in fixed dimension and field
degree was already established by Ahmadi and Parrilo
\cite[Theorem 4.3]{ahmadi-parrilo}. Ahmadi and El Khadir
\cite[Proposition 6.1]{ahmadi-khadir} also compare polynomial and
rational certificate complexity (a degree-five family). Our degree result concerns approach
to the explicit boundary above. Compactness, positive harmonic
representations, and the binary-form sum-of-squares principle are
classical tools, not new general converse theorems.

\begin{remark}\label{rem:inside}
For \(a>0\) one has \(1+a<(1+a^{2/3})^{3/2}\le\sqrt2\,(1+a)<2(1+a)\).
Indeed, for \(x,y>0\) and \(p=2/3\), the \(p\)-power sum satisfies
\(x+y\le(x^p+y^p)^{1/p}\le2^{1/p-1}(x+y)\); take \(x=1,y=a\).
So the new threshold is strictly below the earlier sufficient
nonexistence bound, and the polynomial and analytic obstruction of
\cite{liu} covers only part of the excluded region.
\end{remark}

\section{Polynomial certificates as positive harmonic data}
Use normalized circle averages
\(\avg v=(2\pi)^{-1}\int_0^{2\pi}v(\phi)\,d\phi\).
The circular Hilbert transform \(\mathcal H\) is normalized by
\[
 \mathcal H\cos(k\phi)=\sin(k\phi),\quad
 \mathcal H\sin(k\phi)=-\cos(k\phi),\quad
 \mathcal H1=0,\qquad k\ge1.
\]
On smooth real functions it is skew-adjoint for this average.
The formulas used below can equally be verified termwise from
Fourier series.

\begin{lemma}\label{lem:harmonic}
Let \(P\) be a positive definite real homogeneous polynomial of
degree \(2N\), \(N\ge1\), and set
\(p(\theta)=P(\cos\theta,\sin\theta)\).
There exist \(\alpha_1,\ldots,\alpha_N\in\D\) such that
\[
 h(z)=\frac1N\sum_{j=1}^N\frac{1+\alpha_jz}{1-\alpha_jz},
 \qquad
 \frac{p'(\theta)}{2Np(\theta)}
       =\Ima h(\ee^{2i\theta}).
\]
Here \(h(0)=1\), and \(u=\Rea h\) has a strictly positive minimum
on the closed disk. The weak certificate condition is equivalent to
\begin{equation}\label{eq:weak}
 w(\phi):=a-b\cos\phi-\mathcal Hu(\phi)\ge0.
\end{equation}
\end{lemma}
\begin{proof}
The strict Fej\'er--Riesz factorization writes
\[
 p(\theta)=|C|^2\prod_{j=1}^L
        |1-\alpha_j\ee^{2i\theta}|^2,
 \qquad L\le N,\quad |\alpha_j|<1.
\]
This is the factorization of \cite[Lemma III.2 and proof of
Lemma III.3]{liu}; for completeness here is a direct proof.
The even trigonometric polynomial \(p\) has only the frequencies
\(2j\), \(|j|\le N\), so \(q(\psi)=p(\psi/2)=\sum_{|j|\le L}c_j\ee^{ij\psi}\)
is a positive trigonometric polynomial of actual degree \(L\le N\),
with \(c_{-j}=\bar c_j\) and \(c_L\ne0\). If \(L=0\), \(p\) is constant
and nothing is to be proved. Otherwise
\(A(z)=z^L\sum_jc_jz^j\) is a polynomial of degree \(2L\) with
\(A(0)\ne0\), \(A(\ee^{i\psi})=\ee^{iL\psi}q(\psi)\ne0\), and
\(A(z)=z^{2L}\overline{A(1/\bar z)}\). Hence its roots are not on
the unit circle and come in pairs \(\zeta,1/\bar\zeta\) with equal
multiplicity. Let \(\zeta_1,\dots,\zeta_L\) be the roots in \(\D\)
(none is zero). For \(|z|=1\) one has \(|z-\zeta|=|1-\bar\zeta z|\) and
\(|z-1/\bar\zeta|=|1-\bar\zeta z|/|\zeta|\), so
\(q(\psi)=|A(\ee^{i\psi})|=K\prod_l|1-\alpha_l\ee^{i\psi}|^2\) with
\(\alpha_l=\bar\zeta_l\), \(K=|c_L|/\prod_l|\zeta_l|>0\). Put
\(\psi=2\theta\) and \(|C|^2=K\).
Pad the list with zero parameters to length \(N\).
Each kernel has real part
\[
 \Rea\frac{1+\alpha z}{1-\alpha z}
       =\frac{1-|\alpha z|^2}{|1-\alpha z|^2}>0
       \qquad(|z|\le1).
\]
Differentiating the logarithm of a factor gives twice the imaginary
part of that kernel, proving the angular identity.
The imaginary part of \(h\) is \(\mathcal Hu\) because \(h(0)=1\).
Finally,
\[
 L_fP=2N r^{2N+2}p(\theta)
       \left[-a+b\cos2\theta+
              \frac{p'(\theta)}{2Np(\theta)}\right],
\]
which proves \eqref{eq:weak}.
\end{proof}

Conversely, an equal-weight average of these kernels yields a
polynomial of degree \(2N\). Explicitly, put
\begin{equation}\label{eq:quadratic}
 Q_\alpha(x,y)=(1+|\alpha|^2)(x^2+y^2)
       -2\Rea\alpha(x^2-y^2)+4\Ima\alpha\,xy.
\end{equation}
Its eigenvalues are \((1-|\alpha|)^2,(1+|\alpha|)^2\).
The product \(P=\prod_jQ_{\alpha_j}\) is positive definite and
satisfies the angular identity in Lemma~\ref{lem:harmonic}.
A strict margin in \eqref{eq:weak} therefore gives a strict
polynomial certificate.

\section{The strict obstruction}
Rotation by \(\pi/2\) changes the sign of \(b\), so assume \(b\ge0\).
For \(c>0\), define
\begin{equation}\label{eq:test}
 \begin{gathered}
 s=\sqrt{1+c^2},\qquad \eta=\frac{c+i}{s},\\
 M_c(z)=c^2-1+2ic\frac{1+\eta z}{1-\eta z},
 \qquad f_c(z)=\sqrt{M_c(z)}.
 \end{gathered}
\end{equation}
The Cayley transform has positive real part on \(\D\); hence
\(M_c\) takes values in the upper half-plane. We choose the
analytic square root in the first quadrant. In particular,
\begin{equation}\label{eq:fcoeff}
 f_c(0)=c+i,\qquad f_c'(0)=\frac{2ic}{s}.
\end{equation}

For \(0<\sigma<1\), write
\(f_c(\sigma\ee^{i\phi})=F_\sigma(\phi)+iG_\sigma(\phi)\).
Both functions are positive and smooth. Their means and first
coefficient give
\begin{equation}\label{eq:testmeans}
 \avg{F_\sigma}=c,\quad \avg{G_\sigma}=1,\quad
 \avg{G_\sigma\cos\phi}=\frac{\sigma c}{s},\quad
 F_\sigma+\mathcal HG_\sigma=c.
\end{equation}
If \(u,w\) come from a weak polynomial certificate, then
\(\avg u=1\), \(u\ge\varepsilon>0\), and \(w\ge0\).
Using \eqref{eq:weak} and skew-adjointness,
\begin{align}
 \avg{uF_\sigma+wG_\sigma}
 &=\avg{u(F_\sigma+\mathcal HG_\sigma)}
      +a\avg{G_\sigma}-b\avg{G_\sigma\cos\phi}\nonumber\\
 &=c+a-\frac{b\sigma c}{s}\ge\varepsilon c.
 \label{eq:pairing}
\end{align}
Every function in this computation is smooth for \(\sigma<1\);
no singular boundary pairing is assumed.
Letting \(\sigma\uparrow1\) yields
\begin{equation}\label{eq:dual}
 b\le\frac{(a+c)s}{c}-\varepsilon s
      <\frac{(a+c)\sqrt{1+c^2}}c.
\end{equation}
The derivative of the last expression is
\[
 \frac{c^3-a}{c^2\sqrt{1+c^2}}.
\]
Its minimum occurs at \(c=a^{1/3}\), with value
\((1+a^{2/3})^{3/2}\). This proves the strict necessary
inequality, including nonexistence at equality. The positive
minimum of \(u\), rather than mere nonnegativity, is what
excludes the boundary.

\section{An explicit limiting pair on the boundary}
The converse needs more than a necessary moment inequality.
We construct positive harmonic data at the critical parameters,
then approximate them while retaining a positive margin.

Fix \(c>0\), with \(s,\eta,f_c\) as above, and put
\[
 a_0=c^3,\qquad b_0=s^3.
\]
Define
\begin{align}
 q_c(z)&=1-c^4+\frac{b_0}{2}
          \bigl[(c-i)z+(c+i)z^{-1}\bigr],\label{eq:q}\\
 h_c(z)&=ia_0-\frac{ib_0}{2}(z+z^{-1})
                    +i\frac{q_c(z)}{f_c(z)}.\label{eq:h}
\end{align}

\begin{lemma}\label{lem:critical}
The function \(h_c\) is analytic in \(\D\) (it extends analytically across
zero), continuous on \(\overline\D\), and satisfies \(h_c(0)=1\).
Its boundary data obey
\[
 u_c=\Rea h_c\ge0,\qquad
 w_c=a_0-b_0\cos\phi-\Ima h_c(\ee^{i\phi})\ge0,
\]
and \(u_cw_c\equiv0\): \(w_c\) vanishes on one closed arc of the circle
and \(u_c\) on the complementary closed arc. Moreover \(\Rea h_c>0\) in
\(\D\).
\end{lemma}
\begin{proof}
\emph{Step 1: algebra.} Since \(c+i=s\eta\) and \(c-i=s\bar\eta\),
\begin{equation}\label{eq:Mfactor}
 M_c(z)=\frac{(c+i)^2-\eta z(c-i)^2}{1-\eta z}
        =s^2\frac{\eta^2-\bar\eta z}{1-\eta z},
\end{equation}
because \(\eta\bar\eta^2=\bar\eta\). Using
\(\eta^2+\bar\eta^2=2(c^2-1)/s^2\) and \(1-c^4=s^2(1-c^2)\), expansion of the
right-hand side shows
\begin{equation}\label{eq:qfactor}
 q_c(z)=-\frac{s^2}{2}M_c(z)(1-\eta z)(1-\bar\eta/z)
       =-\frac{s^4}{2}(\eta^2-\bar\eta z)(1-\bar\eta/z).
\end{equation}
So \(M_c\) is a M\"obius map with its only pole at \(z=\bar\eta\) and
its only zero at \(z=\eta^3\) (since \(\eta^2/\bar\eta=\eta^3\)). Both lie
on the unit circle and are distinct, because \(\eta^4\ne1\) for \(c>0\).

\emph{Step 2: branch and interior.} \(M_c\) is analytic near
\(\overline\D\setminus\{\bar\eta\}\), and in \(\D\)
\(\Ima M_c=2c\,\Rea\frac{1+\eta z}{1-\eta z}>0\). The square root with
values in the closed first quadrant is continuous on the closed upper
half-plane and analytic on the open one. Hence \(f_c\) is analytic and
zero-free in \(\D\) and continuous on
\(\overline\D\setminus\{\bar\eta\}\). By \eqref{eq:fcoeff},
\(f_c(0)=c+i\) and \(f_c'(0)=M_c'(0)/(2f_c(0))=2ic/s\).
In \(q_c/f_c\) only the term in \(z^{-1}\) is singular at \(0\). Its
coefficient is \(\frac{b_0}{2}\frac{c+i}{f_c(0)}=\frac{b_0}{2}\), which
cancels the term \(-\frac{ib_0}{2}z^{-1}\) of \eqref{eq:h}. The
constant coefficient of \(h_c\) is
\[
 ia_0+\frac{i(1-c^4)}{f_c(0)}
       -\frac{ib_0}{2}\frac{f_c'(0)}{f_c(0)}
 =ic^3+(1-c^2)(1+ic)+c(c-i)=1 ,
\]
using \(1/(c+i)=(c-i)/s^2\). Hence \(h_c\) is analytic in \(\D\) and
\(h_c(0)=1\).

\emph{Step 3: the circle away from the exceptional points.} Write
\(\eta=\ee^{i\beta}\), \(z=\ee^{i\phi}\), \(\psi=\phi+\beta\). Then
\[
 M_c(\ee^{i\phi})=c^2-1-2c\cot(\psi/2),
\]
which is real and strictly increasing for \(\psi\in(0,2\pi)\), from
\(-\infty\) to \(+\infty\). So \(M_c<0\) on one open arc and \(M_c>0\)
on the complementary open arc, the zero \(\eta^3\) and the pole
\(\bar\eta\) being the endpoints. The boundary root is
\(f_c=\sqrt{M_c}>0\) where \(M_c>0\), and \(f_c=iG\) with
\(G=\sqrt{-M_c}>0\) where \(M_c<0\). On the circle
\(1-\bar\eta/z=\overline{1-\eta z}\), so by \eqref{eq:qfactor}
\(q_c=-\frac{s^2}{2}M_c|1-\eta z|^2\) is real: negative where
\(M_c>0\) and positive where \(M_c<0\). The terms
\(ia_0-ib_0\cos\phi\) of \eqref{eq:h} are purely imaginary there.
Where \(M_c>0\), \(iq_c/f_c\) is purely imaginary, so \(u_c=0\) and
\(w_c=-q_c/f_c>0\). Where \(M_c<0\), \(iq_c/f_c=q_c/G\) is real, so
\(u_c=q_c/G>0\), \(\Ima h_c=a_0-b_0\cos\phi\) and \(w_c=0\).

\emph{Step 4: the two exceptional points.} For
\(z\in\overline\D\setminus\{0,\bar\eta\}\) we have \(|f_c|^2=|M_c|\) and,
by \eqref{eq:Mfactor}, \(|M_c|=s^2|z-\eta^3|/|z-\bar\eta|\) (use
\(|\eta^2-\bar\eta z|=|z-\eta^3|\) and \(|1-\eta z|=|z-\bar\eta|\)).
With \(|1-\bar\eta/z|=|z-\bar\eta|/|z|\), \eqref{eq:qfactor} gives
\[
 \Bigl|\frac{q_c(z)}{f_c(z)}\Bigr|
  =\frac{s^2}{2}|M_c|^{1/2}\,|1-\eta z|\,|1-\bar\eta/z|
  =\frac{s^3}{2}\,
    \frac{|z-\eta^3|^{1/2}\,|z-\bar\eta|^{3/2}}{|z|}.
\]
This is bounded and tends to \(0\) as \(z\to\eta^3\) or
\(z\to\bar\eta\) within \(\overline\D\), with orders \(1/2\) and
\(3/2\) respectively. Since \(q_c/f_c\) is continuous elsewhere on
\(\overline\D\setminus\{0\}\), \(h_c\) extends continuously to the
whole closed disk, with \(h_c=ia_0-ib_0\cos\phi\) at both exceptional
points. There \(u_c=0\) and \(w_c=0\), so the inequalities of Step~3
hold at the endpoints of both arcs.

\emph{Step 5: the interior.} \(\Rea h_c\) is harmonic in \(\D\),
continuous on \(\overline\D\) and \(\ge0\) on the circle. By the minimum
principle \(\Rea h_c\ge0\) in \(\D\), and strictly positive because
\(\Rea h_c(0)=1\).
\end{proof}

This limiting pair is not a polynomial certificate at
\((a_0,b_0)\). Its real part vanishes on an arc, unlike the
strictly positive boundary real part in Lemma~\ref{lem:harmonic}.

\section{Constructing a finite polynomial certificate}
By the rotation by \(\pi/2\) assume \(b\ge0\), and put
\begin{equation}\label{eq:a0}
 a_*(b)=\bigl(\max\{b^{2/3}-1,0\}\bigr)^{3/2},
\end{equation}
so that \(b<(1+a^{2/3})^{3/2}\) is equivalent to \(a>a_*(b)\).
We first fix a \emph{critical function} \(h_*\) for each \(b\ge0\):
an analytic function on \(\D\), continuous on \(\overline\D\), with
\(h_*(0)=1\), \(\Rea h_*\ge0\) on the circle and
\begin{equation}\label{eq:crit}
 \Ima h_*(\ee^{i\phi})+b\cos\phi\le a_*(b)\qquad(\phi\in\R).
\end{equation}
\begin{itemize}
\item If \(b>1\), take \(c=\sqrt{b^{2/3}-1}\), so that
\(b_0=s^3=b\) and \(a_0=c^3=a_*(b)\), and \(h_*=h_c\). Then
\eqref{eq:crit} is the inequality \(w_c\ge0\) of
Lemma~\ref{lem:critical}.
\item If \(0\le b\le1\), take \(h_*(z)=1-ibz\). On the circle
\(\Rea h_*=1+b\sin\phi\ge0\) and \(\Ima h_*+b\cos\phi=0=a_*(b)\).
This treats \(b=0\) and \(b=1\) without the singular endpoint
\(c=0\) of the critical formula.
\end{itemize}
In both cases the measure
\(d\mu(t)=\Rea h_*(\ee^{it})\,dt/2\pi\) has a continuous density, hence
no atoms, and total mass \(\Rea h_*(0)=1\). Poisson's formula for
\(\Rea h_*\) and its harmonic conjugate, together with \(h_*(0)=1\),
give for \(0<\sigma<1\)
\begin{equation}\label{eq:herglotz}
 h_*(\sigma z)=\int_0^{2\pi}K_\sigma(t,z)\,d\mu(t),
 \qquad
 K_\sigma(t,z)=\frac{1+\sigma\ee^{-it}z}{1-\sigma\ee^{-it}z},
 \quad|z|\le1 ;
\end{equation}
the imaginary constant is zero because \(h_*(0)=1\).

\emph{Poisson smoothing.} \(\Ima h_*(\sigma\ee^{i\phi})\) and
\(b\sigma\cos\phi\) are the Poisson integrals at radius \(\sigma\) of
\(\Ima h_*(\ee^{i\phi})\) and \(b\cos\phi\). Integrating
\eqref{eq:crit} against the positive Poisson kernel gives
\(\Ima h_*(\sigma\ee^{i\phi})+b\sigma\cos\phi\le a_*(b)\), hence
\begin{equation}\label{eq:smoothmargin}
 \Ima h_*(\sigma\ee^{i\phi})+b\cos\phi
       \le a_*(b)+b(1-\sigma).
\end{equation}

\emph{Discretization.} Approximating \(\mu\) by equal-weight empirical
measures, with the kernel continuous jointly in \(t\) and
\(|z|\le1\), gives uniform approximation on the closed disk. Thus
for \(\sigma<1\) with \(a_*(b)+b(1-\sigma)<a\), and then for
sufficiently large \(N\),
\[
 h_N(z)=\frac1N\sum_{j=1}^N
        \frac{1+\alpha_jz}{1-\alpha_jz},
 \qquad \alpha_j=\sigma\ee^{-it_j},
\]
retains a strict margin in \eqref{eq:smoothmargin}. The product of the
quadratic forms \eqref{eq:quadratic} is
then a strict certificate. The next proposition gives an explicit
choice of a sufficient \(N\), valid for every \(b\ge0\).

\begin{proposition}[A sufficient degree bound]\label{prop:degree}
Let \(b\ge0\), \(a_*=a_*(b)\) as in \eqref{eq:a0}, and
\(\delta=a-a_*>0\). Set
\[
 \tau=\min\{1/2,\delta/(4b)\}\ \ (\tau=1/2\text{ if }b=0),\qquad
 \sigma=1-\tau,\qquad
 N=\max\left\{1,\left\lceil
                  \frac{16\pi\sigma}{\delta\tau^2}
                    \right\rceil\right\}.
\]
There is a strict polynomial certificate of degree \(2N\)
whose normalized angular margin
\(a-b\cos2\theta-p'/(2Np)\) is at least \(\delta/2\).
If \(0<\delta\le2b\), one may use
\[
 N\le\left\lceil\frac{256\pi b^2}{\delta^3}\right\rceil,
\]
and if \(\delta>2b\) (in particular for \(b=0\)), one may use
\(N\le\max\{1,\lceil32\pi/\delta\rceil\}\).
This is a sufficient bound, not the minimum degree.
\end{proposition}
\begin{proof}
Use \(h_*\) and \(\mu\) as above. Partition \([0,2\pi]\) into \(N\)
successive quantile intervals
\(I_j\), each of \(\mu\)-mass \(1/N\) (possible since \(\mu\) has no
atoms), and choose \(t_j\in I_j\).
Intervals may cross zero-density gaps; their lengths still sum
to \(2\pi\).
For the kernel \(K_\sigma(t,z)\) in \eqref{eq:herglotz}, with
\(w=\sigma\ee^{-it}z\), \(\partial_tK_\sigma=-2iw/(1-w)^2\), so
\[
 \sup_{t,|z|\le1}|\partial_tK_\sigma(t,z)|
       \le\frac{2\sigma}{(1-\sigma)^2}=\frac{2\sigma}{\tau^2}.
\]
The integral-minus-empirical error is therefore at most
\[
 \sum_j\int_{I_j}\frac{2\sigma}{\tau^2}|I_j|\,d\mu
       =\frac{2\sigma}{\tau^2}\cdot\frac1N\sum_j|I_j|
       =\frac{4\pi\sigma}{N\tau^2}\le\frac\delta4.
\]
The smoothing loss in \eqref{eq:smoothmargin} is \(b\tau\le\delta/4\),
so \(\Ima h_N+b\cos\phi\le a_*+\delta/2\), a margin at least
\(\delta/2\) in \(a\). If \(\delta\le2b\), then \(\tau=\delta/(4b)\le1/2\) and
\(\sigma\le1\) give
\(N\le\lceil16\pi\cdot16b^2/\delta^3\rceil\). If \(\delta>2b\), then
\(\tau=\sigma=1/2\) and \(N=\max\{1,\lceil32\pi/\delta\rceil\}\).
\end{proof}

Locating the quantiles with rigorous numerical enclosures is an
additional computational task. Proposition~\ref{prop:degree} is an
analytic construction and degree bound, not a claim that the
floating-point experiment below implements a certified algorithm.

\section{Local analytic certificates and degree divergence}\label{sec:local}
We complete the equivalences of Theorem~\ref{thm:main}.
Let \(V\) be a local real-analytic weak certificate, with Taylor
expansion \(V=P_m+P_{m+1}+\cdots\) at the origin, where \(P_j\) is
homogeneous of degree \(j\) and \(P_m\ne0\).
Along each ray, \(V(tx)=t^mP_m(x)+O(t^{m+1})>0\) gives \(P_m\ge0\).
Since \(f\) is homogeneous of degree three,
\(L_fV(tx)=t^{m+2}L_fP_m(x)+O(t^{m+3})\le0\), so \(L_fP_m\le0\)
(the polynomial \(L_fP_m\) may well vanish identically; only the
inequality is used). With \(P_m(r,\theta)=r^mp_m(\theta)\) and
\eqref{eq:polar}, this reads
\begin{equation}\label{eq:leading}
 p_m'(\theta)+m(-a+b\cos2\theta)p_m(\theta)\le0.
\end{equation}
\begin{lemma}\label{lem:intfactor}
If \(p\ge0\) is a continuous, \(2\pi\)-periodic, piecewise \(C^1\)
function, not identically zero, satisfying
\(p'+m\,g\,p\le0\) for a continuous periodic \(g\), then \(p>0\)
everywhere.
\end{lemma}
\begin{proof}
Let \(\rho(\theta)=\exp\int_0^\theta m\,g\,ds>0\). Then
\((\rho p)'\le0\). If \(p(\theta_0)=0\), the function \(\rho p\) is
nonnegative, nonincreasing on \([\theta_0,\theta_0+2\pi]\) and zero at
\(\theta_0\), hence zero throughout, and by periodicity \(p\equiv0\).
\end{proof}
Hence \(P_m\) is positive definite. Its degree \(m\) is even, since
\(P_m(-x)=(-1)^mP_m(x)\) and \(P_m\ge0\), \(P_m\ne0\). So \(P_m\) is a
positive definite homogeneous weak polynomial certificate, and the strict
obstruction of Section~3 applies to it. This is the leading-term
mechanism already used in \cite[Sec.~III-D]{liu} for the single field
\((a,b)=(1,5)\); the argument is unchanged for every \(a>0\), \(b\in\R\),
so combining it with \cite[Theorem IV.3]{liu} also gives local analytic
nonexistence for \(|b|\ge2(a+1)\).
Conversely, the constructed polynomial is analytic.
Together with the preceding sections this proves
Theorem~\ref{thm:main}.

\begin{proposition}\label{prop:blowup}
For fixed \(b>1\), the minimum degree among positive homogeneous
polynomial weak certificates tends to infinity as
\(a\downarrow a_*(b)=(b^{2/3}-1)^{3/2}\).
\end{proposition}
\begin{proof}
Suppose not. Then there are \(N\) and \(a_n\downarrow a_*\) and positive
definite homogeneous weak certificates \(P_n\) of degree at most
\(2N\) at parameters \((a_n,b)\); passing to a subsequence, all have the
same degree \(2N'\le2N\). Let \(p_n(\theta)=P_n(\cos\theta,\sin\theta)\)
and normalize \(\max p_n=1\) (scaling \(P_n\)). The map from binary forms
of degree \(2N'\) to trigonometric polynomials with frequencies
\(0,\pm2,\dots,\pm2N'\) is a linear isomorphism of spaces of dimension
\(2N'+1\). The coefficients of \(P_n\) are bounded, because
\(\max|p_n|=1\) and all norms on this finite-dimensional space are
equivalent. A subsequence therefore converges, with all derivatives,
to \(p=P(\cos\theta,\sin\theta)\) for a binary form \(P\) of degree
\(2N'\) (possibly of lower Fourier degree, which does not matter),
with \(p\ge0\) and \(\max p=1\). Passing to the limit in
\(p_n'+2N'(-a_n+b\cos2\theta)p_n\le0\) gives the same inequality for
\(p\) at \(a=a_*\). By Lemma~\ref{lem:intfactor}, \(p>0\), so \(P\) is a
positive definite homogeneous polynomial with \(L_fP\le0\) at the
excluded boundary \(b^{2/3}=1+a_*^{2/3}\). This contradicts
Theorem~\ref{thm:main}.
\end{proof}

This compactness method has antecedents in
\cite[Theorem 4.3]{ahmadi-parrilo}.
Propositions~\ref{prop:degree} and \ref{prop:blowup}
do not give matching upper and lower degree rates.

\section{Single harmonics, angular speed and sum of squares}
The same extremal argument applies to a family in every odd degree.
Let \(k\ge1\), \(a>0\), \(b\in\R\), and \(\omega\ne0\), and consider
\begin{equation}\label{eq:higher}
 \dot r=r^{2k+1}(-a+b\cos2k\theta),
 \qquad \dot\theta=\omega r^{2k}.
\end{equation}
It is a polynomial vector field of degree \(2k+1\):
\[
 f(x,y)=
 \bigl[-a(x^2+y^2)^k+b\Rea(x+iy)^{2k}\bigr](x,y)
       +\omega(x^2+y^2)^k(-y,x).
\]

\begin{corollary}\label{cor:higher}
For \eqref{eq:higher}, homogeneous polynomial weak, strict
polynomial, and local real-analytic weak certificates exist
exactly when
\begin{equation}\label{eq:highercriterion}
 |b|^{2/3}<a^{2/3}+|\omega|^{2/3}.
\end{equation}
All these fields are globally asymptotically stable.
The degree bound of Proposition~\ref{prop:degree}, applied to
\(a/|\omega|\) and \(|b|/|\omega|\), becomes \(2kN\) for these fields.
\end{corollary}
\begin{proof}
\emph{Reductions.} Multiplying the field by a positive constant
does not change the certificates, since \(L_{\lambda f}V=\lambda L_fV\);
this replaces \((a,b,\omega)\) by \((a,b,\omega)/|\omega|\). The
reflection \((x,y)\mapsto(x,-y)\) is a linear change of variables
sending \(\theta\mapsto-\theta\); it leaves \(\cos2k\theta\) and
\(r\) unchanged and reverses the sign of \(\theta'=\omega r^{2k}\), so
it maps \eqref{eq:higher} with \(\omega\) to \eqref{eq:higher} with
\(-\omega\), and polynomial or analytic certificates to certificates of
the same kind and degree. The rotation by \(\pi/(2k)\) maps
\(\cos2k\theta\) to \(-\cos2k\theta\) and keeps \(\theta'\), so it
replaces \(b\) by \(-b\). Hence we may assume \(\omega=1\), \(b\ge0\);
\eqref{eq:highercriterion} becomes \(b^{2/3}<1+a^{2/3}\)
(with \(a,b\) replaced by \(a/|\omega|,|b|/|\omega|\)).

\emph{The Lie derivative.} For a homogeneous \(P=r^mp(\theta)\) and
\(\omega=1\),
\[
 L_fP=\dot r\,\partial_rP+\dot\theta\,\partial_\theta P
     =r^{m+2k}\bigl[m(-a+b\cos2k\theta)p+p'\bigr].
\]
Let \(m=2N\) and \(P\) be positive definite. By Lemma~\ref{lem:harmonic},
\(p'/(2Np)=\Ima h(\ee^{2i\theta})\); with \(\phi=2\theta\), \(L_fP\le0\)
is exactly
\begin{equation}\label{eq:weakk}
 w(\phi):=a-b\cos(k\phi)-\mathcal Hu(\phi)\ge0,
 \qquad u=\Rea h,\ \avg u=1,\ u\ge\varepsilon>0 .
\end{equation}

\emph{Necessity (any \(k\), any candidate polynomial).}
Use the test \(f_c(\sigma z^k)\) in place of
\(f_c(\sigma z)\), \(z=\ee^{i\phi}\): write
\(f_c(\sigma\ee^{ik\phi})=F_\sigma(\phi)+iG_\sigma(\phi)\).
This is the boundary value of the function \(z\mapsto f_c(\sigma z^k)\),
analytic in \(\D\) and continuous on \(\overline\D\), so
\(G_\sigma=\mathcal HF_\sigma+\avg{G_\sigma}\) and
\(F_\sigma+\mathcal HG_\sigma=\avg{F_\sigma}\), with
\(\avg{F_\sigma}=\Rea f_c(0)=c\) and
\(\avg{G_\sigma}=\Ima f_c(0)=1\). Because \(F_\sigma+iG_\sigma\) takes
values in the open first quadrant, \(F_\sigma,G_\sigma>0\). Finally,
expanding \(f_c(\zeta)=\sum_{n\ge0}\varphi_n\zeta^n\), only
\(n=1\) contributes to \(\avg{G_\sigma\cos k\phi}\), namely
\(\Ima(\sigma\varphi_1)/2\), and \(\varphi_1=f_c'(0)=2ic/s\):
\[
 \avg{G_\sigma\cos k\phi}=\frac{\sigma c}{s},
 \qquad F_\sigma+\mathcal HG_\sigma=c.
\]
(These are the identities \eqref{eq:testmeans} for \(k=1\); they were
also checked numerically for \(k=1,2,3\), see Section~\ref{sec:verif}.)
The pairing \eqref{eq:pairing} then holds verbatim with
\(\cos\phi\) replaced by \(\cos k\phi\): by skew-adjointness of
\(\mathcal H\),
\(\avg{uF_\sigma+wG_\sigma}=\avg{u(F_\sigma+\mathcal HG_\sigma)}
+a\avg{G_\sigma}-b\avg{G_\sigma\cos k\phi}=c+a-b\sigma c/s\ge\varepsilon c\).
As in Section~3 this gives \(b<(a+c)\sqrt{1+c^2}/c\) for every \(c>0\),
hence \(b<(1+a^{2/3})^{3/2}\). No \(k\)-fold angular symmetry of the
candidate is assumed.

\emph{Sufficiency.} Let \(b\ge0\), \(a>a_*(b)\), and let
\(\alpha_1,\dots,\alpha_N\) be the nodes of Proposition~\ref{prop:degree},
with \(h_N\) the corresponding strict-margin function:
\(\Ima h_N(\ee^{i\psi})+b\cos\psi\le a-\delta/2\) for all \(\psi\). Put
\[
 P(x,y)=\prod_{j=1}^N\bigl|(x-iy)^k-\alpha_j(x+iy)^k\bigr|^2
       =\prod_{j=1}^N\Bigl[(1+|\alpha_j|^2)(x^2+y^2)^k
         -2\Rea\bigl(\alpha_j(x+iy)^{2k}\bigr)\Bigr].
\]
Each factor is a positive definite binary form of degree \(2k\) (on the
circle it equals \(|1-\alpha_j\ee^{2ik\theta}|^2\ge(1-|\alpha_j|)^2\)),
and \(P\) has degree \(2kN\). Differentiating the logarithm of each
factor, with \(\zeta=\ee^{2ik\theta}\), gives
\(\frac{d}{d\theta}\log|1-\alpha\zeta|^2=4k\,\Ima\frac{\alpha\zeta}{1-\alpha\zeta}
=2k\,\Ima\frac{1+\alpha\zeta}{1-\alpha\zeta}\), hence
\[
 \frac{p'(\theta)}{2kNp(\theta)}=\Ima h_N(\ee^{2ik\theta}),\qquad
 \frac{L_fP}{2kN\,r^{2kN+2k}p}\le-\frac\delta2<0
\]
on the circle by the formula for \(L_fP\) above (with \(m=2kN\),
\(\psi=2k\theta\)). Hence \(P\) is a strict polynomial certificate of
degree \(2kN\). The reductions above restore arbitrary \(\omega\ne0\)
and \(b\in\R\) without changing the degree. The reflection and rotation
reductions, the product formula and the sign of \(L_fP\) were also checked
numerically in Cartesian coordinates for several
\((k,a,b,\omega)\), including \(k=2,3\), \(\omega<0\) and \(b<0\)
(Section~\ref{sec:verif}).

\emph{Local analytic certificates.} The argument of Section~\ref{sec:local}
carries over: for a local analytic weak certificate with leading term
\(P_m\) one has \(P_m\ge0\) and
\(L_fP_m=r^{m+2k}[p_m'+m(-a+b\cos2k\theta)p_m]\le0\)
(\(\omega=1\)). Lemma~\ref{lem:intfactor} gives that \(P_m\) is
positive definite, hence of even degree, and the necessity argument above
applies to it.

\emph{Stability.} For \(\omega\ne0\) let
\(H=r^2\exp[-\frac{b}{k\omega}\sin2k\theta]\). Then
\(\partial_rH=2H/r\) and
\(\partial_\theta H=-\frac{2b}{\omega}\cos(2k\theta)H\), so
\[
 L_fH=r^{2k+1}(-a+b\cos2k\theta)\frac{2H}{r}
       -\frac{2b}{\omega}\cos(2k\theta)\,\omega r^{2k}H
      =-2ar^{2k}H.
\]
Its bounded positive angular factor gives \(\ee^{-\beta}r^2\le H\le\ee^{\beta}r^2\),
\(\beta=|b|/(k|\omega|)\), and \(C^1\) regularity at zero (as for
\(k=1\), by \cite[Lemma III.1]{liu}). Decrease of \(H\) bounds all
trajectories and prevents finite-time escape; with \(L_fH<0\) away from
the origin, Lyapunov's direct method gives global asymptotic stability.
\end{proof}

For \(k=1\), write
\(b\cos2\theta+d\sin2\theta=\rho\cos(2\theta-\gamma)\) with
\(\rho=\sqrt{b^2+d^2}\); the rotation of coordinates by \(\gamma/2\)
(which commutes with \(\theta'=\omega r^2\)) turns it into
\(\rho\cos2\theta\). The criterion becomes
\[
 (b^2+d^2)^{1/3}<a^{2/3}+|\omega|^{2/3}.
\]
No result for \(\omega=0\) or a general sum of different
angular harmonics follows from this argument.

\begin{corollary}[Classical binary-form SOS consequence]\label{cor:sos}
The strict certificate region in Theorem~\ref{thm:main} and
Corollary~\ref{cor:higher} is also the region admitting a
homogeneous certificate \(P\) such that \(P\) and \(-L_fP\)
are positive definite sums of squares.
The same sufficient Lyapunov degree bounds apply.
\end{corollary}
\begin{proof}
Any positive definite real binary form \(U\) of degree \(2m\)
is a sum of two squares of real forms of degree \(m\).
Indeed \(U(x,1)\) has positive leading coefficient and no real
roots. Splitting its roots into nonreal conjugate pairs gives
\(U(x,1)=F(x)\overline{F(x)}\). Homogenizing \(F\) yields
\(U=(\Rea F)^2+(\Ima F)^2\).
Apply this classical fact separately to the strict constructed
\(P\) and \(-L_fP\). The reverse implication is immediate.
\end{proof}
Ahmadi and Parrilo \cite[Theorem 5.3]{ahmadi-parrilo} show that any
homogeneous polynomial vector field with a strict homogeneous polynomial
Lyapunov function (\(V\) and \(-\dot V\) positive definite) admits a homogeneous SOS one, \emph{of possibly
higher degree}. In the planar case treated here the classical
two-square fact above keeps the degree; no new general SOS theorem is
claimed.
For any such field the coefficients of \(L_fP\) depend linearly on
those of \(P\). Both \(P\) and \(-L_fP\) are strictly positive on the compact
unit circle, hence (by homogeneity) on \(\R^2\setminus\{0\}\), and
this persists under sufficiently small changes of the coefficient vector
of \(P\). A rational coefficient vector can be chosen in that
neighborhood, so a rational-coefficient strict polynomial certificate
exists as well. This does not assert a two-square decomposition over the
rationals.

\section{Verification and limitations}\label{sec:verif}
The proofs above use explicit analytic identities and inequalities, not
a numerical optimization result. The following computations are
floating-point or high-precision \emph{sanity checks}, not part of any
proof, and none is interval-certified.

\begin{sloppypar}
\emph{Reproducible checks in the supplement.}
The script \texttt{check\_v2.py} (Python with numpy and mpmath, output
\texttt{check\_v2\_output.json}) evaluates, with 40-digit arithmetic,
Lemma~\ref{lem:critical} for \(c\in\{1/8,1/2,1,5/2,8\}\):
\(h_c(0)=1\), \(\Rea h_c\ge0\) on a polar grid of the closed disk
(the smallest value is of rounding size, below \(10^{-38}\) in modulus), \(w_c\ge0\)
and \(u_cw_c=0\) on the circle, identity \eqref{eq:qfactor}, and the
vanishing orders \(0.5\) and \(1.5\) of \(q_c/f_c\) at \(\eta^3\) and
\(\bar\eta\) (measured slopes \(0.500000001\), \(1.500000002\)). It also
builds the products of Corollary~\ref{cor:higher} with \(N=1024\) factors and
\(\sigma=0.98\), directly in Cartesian coordinates, for
\((k,a,b,\omega)\in\{(1,1.2,2,1),(2,1.2,-2,-1),(3,2.4,4,2),(2,0.6,0.8,1),
(2,0.5,0,1)\}\); the largest sampled value of \(L_fP/(2kN\,P)\) on the
circle (20001 angles) is negative in each case (for example \(-0.688\) for the first
two, matching \(a-\max(\Ima h_N+b\cos)\)). Finally it confirms the
identities \(\avg{F_\sigma}=c\), \(\avg{G_\sigma}=1\),
\(\avg{G_\sigma\cos k\phi}=\sigma c/s\), \(F_\sigma+\mathcal HG_\sigma=c\)
for \(k=1,2,3\) to about \(10^{-13}\).
\end{sloppypar}

\emph{An anecdotal example from version 1.} One finite product of
2048 quadratic factors was evaluated at 8192 angular samples for
\(a=1.2\), \(b=2\sqrt2\), \(\sigma=0.98\) (script
\texttt{check\_phase30.py}, output \texttt{phase-checks-30.json}); the
sampled normalized margin was about \(0.1545\). It is illustrative only:
it is not a certified positivity check, not a minimum-degree result,
and nothing in the paper depends on it. A further exploratory
Toeplitz-relaxation program of version 1 is dropped from the text;
its script and stored output remain in the supplement, with their
solver-status caveats, and no conclusion here relies on them.

\emph{Review status.} The manuscript was written with AI assistance.
Version 1 received separate automated reviews by Claude Sonnet 5.5 and
Claude Haiku 5.5; version 2 responds to them. These are not human,
blind or independent peer reviews (the reviewer and the drafting
assistant belong to the same model family), and not formal verification.
No Lean formalization of this argument exists; the Lean development
mentioned in \cite{liu} concerns their Theorem~II.2 only. A rigorously
enclosed implementation of the quantile selection is not supplied; the
quantitative construction is analytic.

\emph{Sources.} The comparison with Liu--Fitzsimmons uses the complete
text of arXiv:2607.16171v1 (the only version found on 21 October 2026).
The cited theorem statements of Ahmadi--Parrilo (Theorems 4.3 and 5.3) and
of Ahmadi--El Khadir (Proposition 6.1) were checked against their
texts, but those works were not read in full. Targeted searches did not locate this
boundary; worldwide priority is unconfirmed, particularly for the
associated harmonic extremal problem. Minimum-degree asymptotics and
general multiharmonic classification remain open here.

\small
\begin{thebibliography}{9}
\bibitem{liu}
J. Liu and M. Fitzsimmons,
\emph{A globally asymptotically stable planar homogeneous polynomial
vector field with no polynomial Lyapunov function},
arXiv:2607.16171v1, 17 July 2026.
\url{https://arxiv.org/abs/2607.16171v1}.
\bibitem{ahmadi-parrilo}
A. A. Ahmadi and P. A. Parrilo,
\emph{Stability of polynomial differential equations:
complexity and converse Lyapunov questions}, arXiv:1308.6833.
Theorem numbers (4.3, 5.3) as in arXiv v1 and the author-hosted manuscript:
\url{https://web.mit.edu/~a_a_a/Public/Publications/Selected_Papers/dTAC.pdf}.
\bibitem{ahmadi-khadir}
A. A. Ahmadi and B. El Khadir,
\emph{On algebraic proofs of stability for homogeneous vector fields},
IEEE Trans. Automat. Control \textbf{65} (2020), 325--332.
Proposition 6.1 and Lemma 6.2 as numbered in the 2018 preprint
(cited as Proposition VI.1 in \cite{liu}):
\url{https://optimization-online.org/wp-content/uploads/2018/03/6502.pdf}.
\end{thebibliography}
\end{document}
